\section{事件响应与复盘}
\subsection{故障归因流程}
每次 Agent 失败后团队会立即启动复盘：首先提取 `failure matrix` 包含命令、触发文件、执行环境及返回码；然后比对当前策略与历史策略，判断是 prompt、工具链还是运行环境的回归。

\begin{table}[H]
\centering
\begin{tabular}{p{0.3\textwidth} p{0.3\textwidth} p{0.3\textwidth}}
\toprule
分类 & 典型触发条件 & 复盘动作 \\
\midrule
Prompt 失效 & 生成与上下文不符 & 把 prompt 重新模板化，并加入新的 examples \\
\midrule
工具链错误 & 命令/测试挂 & 检查 docker 镜像、依赖版本、权限 \\
\midrule
评测环境差异 & 本地通过但 CI 失败 & 对比环境变量和数据版本，做 dev vs prod tests \\
\bottomrule
\end{tabular}
\caption{故障分类与复盘动作。}
\end{table}

\subsection{复盘纪要传播}
复盘结果会以 Markdown 形式写入团队 wiki，内容包括时间、负责人、复盘结论和后续动作，并通过 Slack 频道 \texttt{\#agent-postmortem} 发送提醒。每次复盘还会带上关键截图（如测试 log、 diff）让未来快速定位。

\begin{warningbox}{不要让复盘沉没}
复盘纪要必须在 24 小时内发布，过久会导致场景遗忘。
\end{warningbox}

\begin{knowledgebox}{复盘价值}
复盘不仅是找错，更是把 Agent 的失败样本变成训练 prompt 的课件，把“为什么失败”转化成可重复的规律。
\end{knowledgebox}

\subsection{本章小结}
\begin{itemize}
    \item 快速归因与复盘可防止同类问题反复出现。
    \item 把复盘结果分享给团队可提升整体 Agent 工程成熟度。
\end{itemize}

